\section{Local convergence of random fields}

Throughout, $(E, \mathcal{E})$ is a measurable space, $S$ a set of sites,
$(\Cfg, \mathcal{F}) = (E^S, \mathcal{E}^S)$ the configuration space, and
$\Prob(\Cfg, \mathcal{F})$ the set of probability measures on it.

\subsection*{The algebra of local events}

\begin{definition}[Local events $\mathcal{F}^0$, Georgii (4.1)]
    \label{def:lc-local-events}
    \uses{def:cylinder-event}
    \lean{MeasureTheory.localEvents}
    \leanok

    The \textbf{algebra of local events} is
    \[
        \mathcal{F}^0 \;=\; \bigcup_{\Lambda \in \Fin} \mathcal{F}_\Lambda,
    \]
    the measurable cylinder events of $\Cfg = E^S$.
\end{definition}

\begin{lemma}[Two descriptions of $\mathcal{F}^0$]
    \label{lem:lc-mem-local-events}
    \uses{def:lc-local-events, def:cylinder-event}
    \lean{MeasureTheory.mem_localEvents_iff_cylinderEvents,
      MeasureTheory.mem_localEvents_iff_exists_finsetRestrict_preimage,
      MeasureTheory.mem_localEvents_of_cylinderEvents,
      MeasureTheory.finsetRestrict_preimage_mem_localEvents}
    \leanok

    For $A \subseteq \Cfg$ the following are equivalent:
    \begin{enumerate}
        \item $A \in \mathcal{F}^0$;
        \item $A \in \mathcal{F}_\Lambda$ for some $\Lambda \in \Fin$;
        \item $A = \pi_\Lambda^{-1}(B)$ for some $\Lambda \in \Fin$ and some measurable
        $B \subseteq E^\Lambda$, where $\pi_\Lambda : E^S \to E^\Lambda$ is the restriction map.
    \end{enumerate}
\end{lemma}
\begin{proof}
    \leanok

    (ii)$\Leftrightarrow$(iii): $\mathcal{F}_\Lambda$ is the $\sigma$-algebra pulled back along
    $\pi_\Lambda$. (i)$\Leftrightarrow$(iii): the membership criterion for measurable cylinders.
\end{proof}

\subsection*{The topology of local convergence}

\begin{definition}[Topology of local convergence, Georgii (4.2)]
    \label{def:lc-topology}
    \uses{def:lc-local-events}
    \lean{MeasureTheory.WithLocalConvergence, MeasureTheory.WithSetwiseTopology}
    \leanok

    The \textbf{topology of local convergence} ($\mathcal{L}$-topology) on
    $\Prob(\Cfg, \mathcal{F})$ is the coarsest topology making every evaluation
    $\nu \mapsto \nu(A)$, $A \in \mathcal{F}^0$, continuous; equivalently, the sets
    \[
        \Big\{ \nu \in \Prob(\Cfg, \mathcal{F}) : \max_{1 \le k \le n} |\nu(A_k) - \mu(A_k)|
        < \varepsilon \Big\},
        \qquad A_1, \dots, A_n \in \mathcal{F}^0,\ \varepsilon > 0,
    \]
    form a neighbourhood base at $\mu$. The same definition applies to an arbitrary family
    $\mathcal{C}$ of measurable sets in place of $\mathcal{F}^0$, and to measures and finite
    measures in place of probability measures.
\end{definition}

\begin{lemma}[Convergence in the $\mathcal{L}$-topology, Georgii (4.2)]
    \label{lem:lc-tendsto-iff}
    \uses{def:lc-topology, def:lc-local-events}
    \lean{MeasureTheory.tendsto_withLocalConvergence_iff,
      MeasureTheory.WithSetwiseTopology.tendsto_prob_iff}
    \leanok

    Let $l$ be a filter on an index set $\iota$. A family $(\mu_i)_{i \in \iota}$ in
    $\Prob(\Cfg, \mathcal{F})$ converges locally to $\mu$ along $l$ if and only if
    $\lim_l \mu_i(A) = \mu(A)$ for every $A \in \mathcal{F}^0$.
\end{lemma}
\begin{proof}
    \leanok

    Convergence in an initial topology induced by a map into a product is convergence of each
    coordinate of the image.
\end{proof}

\begin{lemma}[Evaluation on a local event is continuous]
    \label{lem:lc-continuous-eval}
    \uses{def:lc-topology, def:lc-local-events}
    \lean{MeasureTheory.WithSetwiseTopology.continuous_apply_enn,
      MeasureTheory.WithSetwiseTopology.continuous_apply_real}
    \leanok

    For $A \in \mathcal{F}^0$, both $\mu \mapsto \mu(A) \in \overline{\R}_{\ge 0}$ and
    $\mu \mapsto \mu(A) \in \R$ are continuous on $\Prob(\Cfg, \mathcal{F})$.
\end{lemma}
\begin{proof}
    \leanok

    The first is a coordinate of the inducing evaluation map. For the second,
    $\mu(A) \le 1 < \infty$ and $\overline{\R}_{\ge 0} \to \R$ is continuous away from $\infty$.
\end{proof}

\begin{lemma}[Local events separate probability measures]
    \label{lem:lc-separates}
    \uses{def:lc-local-events}
    \lean{MeasureTheory.separatesOn_localEvents,
      MeasureTheory.WithSetwiseTopology.separatesOn_of_isPiSystem_of_generateFrom}
    \leanok

    If $\mu, \nu \in \Prob(\Cfg, \mathcal{F})$ satisfy $\mu(A) = \nu(A)$ for all
    $A \in \mathcal{F}^0$, then $\mu = \nu$.
\end{lemma}
\begin{proof}
    \leanok

    $\mathcal{F}^0$ is a $\pi$-system generating $\mathcal{F}$, and both measures have total
    mass $1$; apply the uniqueness part of the extension theorem.
\end{proof}

\begin{proposition}[Georgii (4.3)(1): the $\mathcal{L}$-topology is Hausdorff]
    \label{prop:lc-hausdorff}
    \uses{lem:lc-separates, def:lc-topology}
    \lean{MeasureTheory.instT2SpaceWithLocalConvergence,
      MeasureTheory.WithSetwiseTopology.t2Space_probabilityMeasure}
    \leanok

    $\Prob(\Cfg, \mathcal{F})$ with the $\mathcal{L}$-topology is Hausdorff.
\end{proposition}
\begin{proof}
    \leanok

    By Lemma~\ref{lem:lc-separates} the evaluation map
    $\mu \mapsto (\mu(A))_{A \in \mathcal{F}^0}$ is injective, and it is inducing by
    construction; an inducing injection into the Hausdorff product
    $\overline{\R}_{\ge 0}^{\,\mathcal{F}^0}$ is an embedding.
\end{proof}

\begin{remark}[Georgii (4.3)(1): the $\mathcal{L}$-topology is an initial topology]
    \label{rmk:lc-initial-topology}
    \uses{def:lc-topology}
    \lean{MeasureTheory.WithSetwiseTopology.isInducing_evalProb,
      MeasureTheory.WithSetwiseTopology.evalProb}
    \leanok

    The restriction map $p : \mu \mapsto \mu | \mathcal{F}^0$, viewed as the evaluation map into
    $\overline{\R}_{\ge 0}^{\,\mathcal{F}^0}$ with the product topology, is inducing: the
    $\mathcal{L}$-topology is the initial topology of $p$. With
    Proposition~\ref{prop:lc-hausdorff}, $p$ is a homeomorphism onto its image.
\end{remark}

\begin{remark}[Georgii (4.3)(1): $p$ is a bijection]
    \label{rmk:lc-caratheodory}
    \uses{rmk:lc-initial-topology}

    By Carath\'eodory's extension theorem, $p$ is a bijection from $\Prob(\Cfg, \mathcal{F})$ onto
    the set $\Prob(\Cfg, \mathcal{F}^0)$ of $\sigma$-additive normalized nonnegative set functions
    on the algebra $\mathcal{F}^0$.
\end{remark}

\subsection*{Georgii (4.3)(2): the $\mathcal{L}$-topology is the topology of local observables}

\begin{definition}[Local and quasilocal observables, Georgii (2.20)]
    \label{def:lc-local-functions}
    \uses{def:cylinder-event}
    \lean{MeasureTheory.GibbsMeasure.localFunctionsOn,
      MeasureTheory.GibbsMeasure.localFunctions,
      MeasureTheory.GibbsMeasure.quasilocalFunctions}
    \leanok

    $\mathcal{L}_\Lambda$ is the subalgebra of $\ell^\infty(\Cfg; \R)$ of bounded
    $\mathcal{F}_\Lambda$-measurable functions;
    $\mathcal{L} = \bigcup_{\Lambda \in \Fin} \mathcal{L}_\Lambda$, the supremum of the directed
    family $(\mathcal{L}_\Lambda)_{\Lambda \in \Fin}$, is the algebra of bounded local observables;
    $\overline{\mathcal{L}}$, the closure of $\mathcal{L}$ in the supremum norm, is the algebra of
    quasilocal observables.
\end{definition}

\begin{definition}[$\mathcal{L}$-continuous observables]
    \label{def:lc-lcontinuous}
    \uses{def:lc-topology}
    \lean{MeasureTheory.GibbsMeasure.LContinuous}
    \leanok

    A bounded observable $f$ is \textbf{$\mathcal{L}$-continuous} if
    $\mu \mapsto \mu(f) = \int f \, d\mu$ is continuous on $\Prob(\Cfg, \mathcal{F})$ for the
    $\mathcal{L}$-topology.
\end{definition}

\begin{lemma}[The $\mathcal{L}$-continuous observables are uniformly closed]
    \label{lem:lc-isclosed-lcontinuous}
    \uses{def:lc-lcontinuous}
    \lean{MeasureTheory.GibbsMeasure.isClosed_lContinuous}
    \leanok

    The set of bounded measurable $\mathcal{L}$-continuous observables is closed in
    $\ell^\infty(\Cfg; \R)$.
\end{lemma}
\begin{proof}
    \leanok

    If $f_n \to g$ uniformly then $|\nu(g) - \nu(f_n)| \le \|g - f_n\|_\infty$ for every
    probability measure $\nu$, so $\nu \mapsto \nu(f_n)$ converges uniformly to
    $\nu \mapsto \nu(g)$, and a uniform limit of continuous functions is continuous. Boundedness
    and measurability pass to the limit.
\end{proof}

\begin{lemma}[Finite-volume simple observables are $\mathcal{L}$-continuous]
    \label{lem:lc-simple}
    \uses{def:lc-lcontinuous, lem:lc-continuous-eval, lem:lc-mem-local-events}
    \lean{MeasureTheory.GibbsMeasure.lContinuous_of_mem_simpleFunctions}
    \leanok

    If $\Lambda \in \Fin$ and $g = \sum_{k=1}^n a_k \mathbf{1}_{A_k}$ is a simple
    $\mathcal{F}_\Lambda$-measurable function, then $g$ is $\mathcal{L}$-continuous.
\end{lemma}
\begin{proof}
    \leanok

    $\nu(g) = \sum_{y \in \mathrm{ran}\, g} y \cdot \nu(g^{-1}\{y\})$, a finite sum. Each fibre
    $g^{-1}\{y\}$ lies in $\mathcal{F}_\Lambda \subseteq \mathcal{F}^0$, so each summand is
    continuous by Lemma~\ref{lem:lc-continuous-eval}.
\end{proof}

\begin{proposition}[Quasilocal observables have continuous integral]
    \label{prop:lc-quasilocal-continuous}
    \uses{lem:lc-isclosed-lcontinuous, lem:lc-simple, def:lc-local-functions, def:lc-lcontinuous}
    \lean{MeasureTheory.GibbsMeasure.lContinuous_of_mem_quasilocalFunctions}
    \leanok

    Every $f \in \overline{\mathcal{L}}$ is $\mathcal{L}$-continuous.
\end{proposition}
\begin{proof}
    \leanok

    For fixed $\Lambda$, $\mathcal{L}_\Lambda$ is the uniform closure of the simple
    $\mathcal{F}_\Lambda$-measurable functions, which are $\mathcal{L}$-continuous by
    Lemma~\ref{lem:lc-simple}; by Lemma~\ref{lem:lc-isclosed-lcontinuous}, so is
    $\mathcal{L}_\Lambda$, hence $\mathcal{L}$, hence (Lemma~\ref{lem:lc-isclosed-lcontinuous}
    again) $\overline{\mathcal{L}}$.
\end{proof}

\begin{definition}[Indicator of a set as a bounded observable]
    \label{def:lc-indicator}
    \lean{MeasureTheory.GibbsMeasure.indicatorLp, MeasureTheory.GibbsMeasure.coeFn_indicatorLp}
    \leanok

    For $A \subseteq \Cfg$, $\mathbf{1}_A$ regarded as an element of $\ell^\infty(\Cfg; \R)$.
\end{definition}

\begin{lemma}[Indicators of local events are local observables]
    \label{lem:lc-indicator-local}
    \uses{def:lc-indicator, def:lc-local-events, def:lc-local-functions}
    \lean{MeasureTheory.GibbsMeasure.indicatorLp_mem_localFunctions}
    \leanok

    If $A \in \mathcal{F}^0$ then $\mathbf{1}_A \in \mathcal{L}$.
\end{lemma}
\begin{proof}
    \leanok

    By Lemma~\ref{lem:lc-mem-local-events}, $A \in \mathcal{F}_\Lambda$ for some
    $\Lambda \in \Fin$, so $\mathbf{1}_A$ is bounded and $\mathcal{F}_\Lambda$-measurable.
\end{proof}

\begin{theorem}[Georgii (4.3)(2)]
    \label{thm:lc-tendsto-iff-localFunctions}
    \uses{prop:lc-quasilocal-continuous, lem:lc-indicator-local, lem:lc-tendsto-iff, def:lc-local-functions, def:lc-topology}
    \lean{MeasureTheory.GibbsMeasure.tendsto_iff_forall_localFunctions}
    \leanok

    For any filter $l$ on an index set $\iota$, a family $(\mu_i)_{i \in \iota}$ converges to
    $\mu$ in the $\mathcal{L}$-topology along $l$ if and only if $\lim_l \mu_i(f) = \mu(f)$ for
    every $f \in \mathcal{L}$. Hence the $\mathcal{L}$-topology is the coarsest topology for which
    $\nu \mapsto \nu(f)$ is continuous for every $f \in \mathcal{L}$.
\end{theorem}
\begin{proof}
    \leanok

    ($\Rightarrow$) $\mathcal{L} \subseteq \overline{\mathcal{L}}$, so $\nu \mapsto \nu(f)$ is
    continuous by Proposition~\ref{prop:lc-quasilocal-continuous}.
    ($\Leftarrow$) For $A \in \mathcal{F}^0$ take $f = \mathbf{1}_A$
    (Lemma~\ref{lem:lc-indicator-local}); $\nu(\mathbf{1}_A) = \nu(A)$ is finite, so convergence
    of the reals $\mu_i(A)$ gives convergence in $\overline{\R}_{\ge 0}$, which is local
    convergence by Lemma~\ref{lem:lc-tendsto-iff}.
\end{proof}

\subsection*{Georgii (4.3)(3): metrizability for a finite state space}

\begin{lemma}[Countability of $\mathcal{F}^0$]
    \label{lem:lc-countable-local-events}
    \uses{def:lc-local-events, lem:lc-mem-local-events}
    \lean{MeasureTheory.countable_localEvents,
      MeasureTheory.instCountableElemSetForallLocalEventsOfFinite}
    \leanok

    If $S$ is countable and $E$ is finite, then $\mathcal{F}^0$ is countable.
\end{lemma}
\begin{proof}
    \leanok

    By Lemma~\ref{lem:lc-mem-local-events} every local event is $\pi_\Lambda^{-1}(B)$ with
    $\Lambda \in \Fin$ and $B \subseteq E^\Lambda$; there are countably many $\Lambda$ and, for
    each, finitely many $B$.
\end{proof}

\begin{proposition}[Georgii (4.3)(3): metrizability]
    \label{prop:lc-metrizable}
    \uses{lem:lc-countable-local-events, rmk:lc-initial-topology, prop:lc-hausdorff, def:lc-topology}
    \lean{MeasureTheory.instPseudoMetrizableSpaceWithLocalConvergenceOfCountableOfFinite,
      MeasureTheory.instMetrizableSpaceWithLocalConvergenceOfCountableOfFinite}
    \leanok

    If $S$ is countable and $E$ is finite, then the $\mathcal{L}$-topology on
    $\Prob(\Cfg, \mathcal{F})$ is metrizable.
\end{proposition}
\begin{proof}
    \leanok

    By Remark~\ref{rmk:lc-initial-topology} the space is induced from the product
    $\overline{\R}_{\ge 0}^{\,\mathcal{F}^0}$, whose index set is countable by
    Lemma~\ref{lem:lc-countable-local-events}; a countable product of metrizable spaces is
    metrizable, and an inducing map into a pseudometrizable space makes the source
    pseudometrizable. Hausdorff (Proposition~\ref{prop:lc-hausdorff}) and pseudometrizable give
    metrizable.
\end{proof}

\begin{remark}[Georgii (4.3)(3): metrizability is an equivalence]
    \label{rmk:lc-metrizability-converse}
    \uses{prop:lc-metrizable}

    The $\mathcal{L}$-topology is metrizable if and only if $\mathcal{F}^0$ is countable, which
    happens if and only if $E$ is finite. In that case a compatible metric is
    \[
        d(\mu, \nu) \;=\; \sum_{n \ge 1} 2^{-n} \sum_{\zeta \in E^{\Lambda(n)}}
        |\mu(\sigma_{\Lambda(n)} = \zeta) - \nu(\sigma_{\Lambda(n)} = \zeta)|
    \]
    for a cofinal sequence $(\Lambda(n))_{n \ge 1}$ in $\Fin$.
\end{remark}

\begin{corollary}[Random fields over a finite state space form a Polish space]
    \label{cor:lc-polish}
    \uses{prop:lc-metrizable, ex:cp-411-2}
    \lean{MeasureTheory.compactSpace_and_metrizableSpace_withLocalConvergence,
      MeasureTheory.instIsCompletelyMetrizableSpaceWithLocalConvergenceOfCountableOfFiniteOfMeasurableSingletonClass,
      MeasureTheory.instPolishSpaceWithLocalConvergenceOfCountableOfFiniteOfMeasurableSingletonClass}
    \leanok

    If $S$ is countable and $E$ is finite with all singletons measurable, then
    $\Prob(\Cfg, \mathcal{F})$ with the $\mathcal{L}$-topology is compact metrizable, hence
    completely metrizable, hence Polish.
\end{corollary}
\begin{proof}
    \leanok

    Compactness is Example~\ref{ex:cp-411-2}; metrizability is
    Proposition~\ref{prop:lc-metrizable}. A compact metrizable space is complete for every
    compatible metric and separable, hence Polish.
\end{proof}

\begin{remark}[Comparison with the weak topology, Georgii (4.3)(3)]
    \label{rmk:lc-weak-topology}
    \uses{thm:lc-tendsto-iff-localFunctions}

    If $E$ is a metric space and $\Cfg$ carries the product topology, the $\mathcal{L}$-topology
    is finer than the weak topology; the two coincide when $E$ is countable and discrete, every
    local function then being continuous.
\end{remark}

\subsection*{Cylinder functions on the configuration space}

\begin{definition}[Cylinder function]
    \label{def:lc-cylinder-function}
    \lean{MeasureTheory.GibbsMeasure.ConfigurationSpace.IsCylinderFunction}
    \leanok

    A function $f : \Cfg \to F$, with $F$ an arbitrary codomain and no boundedness or
    measurability assumed, is a \textbf{cylinder function} if there is $\Lambda \in \Fin$ with
    $f(\zeta) = f(\eta)$ whenever $\zeta$ and $\eta$ agree on $\Lambda$. Every
    $f \in \mathcal{L}$ (Definition~\ref{def:lc-local-functions}) is a cylinder function.
\end{definition}

\begin{lemma}[Cylinder functions are stable under the algebraic operations]
    \label{lem:lc-cylinder-function-algebra}
    \uses{def:lc-cylinder-function}
    \lean{MeasureTheory.GibbsMeasure.ConfigurationSpace.IsCylinderFunction.const,
      MeasureTheory.GibbsMeasure.ConfigurationSpace.IsCylinderFunction.comp,
      MeasureTheory.GibbsMeasure.ConfigurationSpace.IsCylinderFunction.comp₂,
      MeasureTheory.GibbsMeasure.ConfigurationSpace.IsCylinderFunction.add,
      MeasureTheory.GibbsMeasure.ConfigurationSpace.IsCylinderFunction.mul,
      MeasureTheory.GibbsMeasure.ConfigurationSpace.IsCylinderFunction.smul,
      MeasureTheory.GibbsMeasure.ConfigurationSpace.IsCylinderFunction.neg,
      MeasureTheory.GibbsMeasure.ConfigurationSpace.IsCylinderFunction.sub}
    \leanok

    Constants are cylinder functions; post-composition with any map preserves the property; if
    $f, g$ are cylinder functions so is $F(f, g)$ for any binary $F$, in particular $f + g$,
    $f \cdot g$, $c \cdot f$, $-f$ and $f - g$.
\end{lemma}
\begin{proof}
    \leanok

    A constant depends on $\emptyset$; post-composition preserves the dependency set; two
    functions depending on $\Lambda_f$ and $\Lambda_g$ jointly depend on
    $\Lambda_f \cup \Lambda_g$.
\end{proof}

\begin{lemma}[Product measurable structure is the Borel structure]
    \label{lem:lc-borel}
    \lean{MeasureTheory.GibbsMeasure.ConfigurationSpace.measurableSpace_pi_eq_borel}
    \leanok

    If $S$ is countable and $E$ is a second-countable Borel space, then the product
    $\sigma$-algebra $\mathcal{E}^S$ on $\Cfg = E^S$ is the Borel $\sigma$-algebra of the product
    topology.
\end{lemma}
\begin{proof}
    \leanok

    A countable product of second-countable Borel spaces is a Borel space.
\end{proof}
